% ============================================================
% The Hubble Tension as Heterogeneity Gap:
% Six Tier-2 Theorems in the Generative-Collapse Kernel
%
% A perspective piece framed through RCFT — Recursive Collapse
% Field Theory as diagnostic lens on cosmological measurement
% incoherence.
%
% Compile: pdflatex → bibtex → pdflatex → pdflatex
% ============================================================

\documentclass[
  aps,
  prd,
  twocolumn,
  superscriptaddress,
  nofootinbib,
  floatfix
]{revtex4-2}

% ── packages ──
\usepackage{amsmath,amssymb,amsthm,mathtools}
\usepackage{graphicx}
\usepackage{xcolor}
\definecolor{linkblue}{HTML}{337799}
\usepackage[colorlinks=true,linkcolor=linkblue,citecolor=linkblue,urlcolor=linkblue]{hyperref}
\usepackage{bm}
\usepackage{booktabs}
\usepackage{multirow}
\usepackage{xfrac}
\usepackage{enumitem}
\usepackage{array}
\newcolumntype{P}[1]{>{\raggedright\arraybackslash}p{#1}}

% ── theorem environments ──
\newtheorem*{axiom}{The Axiom}
\newtheorem{theorem}{Theorem}
\newtheorem{definition}[theorem]{Definition}
\newtheorem{lemma}[theorem]{Lemma}
\newtheorem{corollary}[theorem]{Corollary}
\newtheorem{proposition}[theorem]{Proposition}
\newtheorem{identity}[theorem]{Identity}
\newtheorem{remark}{Remark}

% ── UMCP / GCD / RCFT macros ──
\newcommand{\tR}{\tau_{\!R}}
\newcommand{\tRS}{\tau_{\!R}^{*}}
\newcommand{\Gam}{\Gamma}
\newcommand{\eps}{\varepsilon}
\newcommand{\IR}{\infty_{\mathrm{rec}}}
\newcommand{\tolseam}{\mathrm{tol}_{\mathrm{seam}}}
\newcommand{\dd}{\mathrm{d}}
\newcommand{\beq}{\begin{equation}}
\newcommand{\eeq}{\end{equation}}
\newcommand{\INF}{\texttt{INF\_REC}}
\newcommand{\regime}[1]{\textrm{\upshape\scshape#1}}
\newcommand{\conform}{\regime{conformant}}
\newcommand{\nonconform}{\regime{nonconformant}}
\newcommand{\tvec}{\mathbf{c}}
\newcommand{\wvec}{\mathbf{w}}
\newcommand{\IC}{\mathrm{IC}}
\newcommand{\gF}{g_{\!F}}
\newcommand{\hetgap}{\Delta}
\newcommand{\Fid}{F}
\newcommand{\Drft}{\omega}
\newcommand{\Ent}{S}
\newcommand{\Curv}{C}
\newcommand{\Logint}{\kappa}
\newcommand{\lat}[1]{\textit{#1}}

\begin{document}

% ============================================================
% TITLE
% ============================================================
\title{The Hubble Tension as Heterogeneity Gap:\\
Six Tier-2 Theorems in the Generative-Collapse Kernel}

\author{Clement Paulus}
\email{clementpaulus9@gmail.com}
\affiliation{UMCP / GCD / RCFT Canon.
  Repository:
  \href{https://github.com/calebpruett927/GENERATIVE-COLLAPSE-DYNAMICS}{GENERATIVE-COLLAPSE-DYNAMICS}.}

\date{\today}

% ============================================================
\begin{abstract}
We apply the six-metric kernel of Generative Collapse Dynamics
(GCD), governed by the Universal Measurement Contract Protocol
(UMCP), to the Hubble tension --- the persistent $\sim\!5\sigma$
discrepancy between early-universe and late-universe
determinations of the Hubble constant~$H_0$.
Encoding 12 measurements from five methodological classes
(CMB, BAO, distance ladder, time-delay cosmography, stellar
ages, and gravitationally lensed supernovae) as 8-channel
trace vectors $\tvec \in [\eps, 1{-}\eps]^{8}$, we prove
six Tier-2 theorems (T-HT-1 through T-HT-6) with 93 automated
tests and 0 failures.

The central finding is a reinterpretation: the standard framing
``CMB says $H_0 \approx 67$, distance ladder says $H_0 \approx 73$''
is structurally incomplete.
The GCD kernel reveals that late-universe measurements have a
heterogeneity gap $\hetgap_{\mathrm{late}} = 0.254$ that is
$6\times$ larger than $\hetgap_{\mathrm{early}} = 0.042$.
The distance ladder methods (SH0ES, TRGB, JAGB, TDCOSMO)
disagree with \emph{each other} far more than the CMB-based
methods disagree among themselves.
The Hubble tension is not merely ``early vs.\ late'' --- it is
a structural incoherence concentrated in late-universe methodologies,
visible as geometric slaughter in the kernel's multiplicative
integrity~$\IC$.

ACT DR6~\cite{louis2025act_dr6,calabrese2025act_ext} tested
12 extended $\Lambda$CDM models and found none favored over
the standard model --- no Tier-2 closure eliminates the gap.
The JWST VENUS program~\cite{venus2026} offers a resolution
pathway through gravitationally lensed supernovae that are
structurally independent of the Cepheid/TRGB/JAGB calibration
chain.
All results derive from Tier-1 identities
($\Fid + \Drft = 1$, $\IC \leq \Fid$, $\IC = e^{\Logint}$)
applied to physically motivated trace vectors, with duality
verified to machine precision across all 12 measurements.
\end{abstract}

\maketitle

% ============================================================
\section{Introduction}\label{sec:intro}
% ============================================================

The Hubble constant $H_0$ parametrizes the present-day expansion
rate of the universe.
Its value, determined independently from the cosmic microwave
background (CMB) radiation at redshift $z \sim 1100$ and from
calibrated distance ladders in the local universe at
$z \lesssim 0.15$, has yielded two persistently discrepant
clusters of measurements separated by approximately
$5.6\,\mathrm{km\,s^{-1}\,Mpc^{-1}}$
(Planck $67.4 \pm 0.5$~\cite{planck2018vi} versus SH0ES
$73.04 \pm 1.04$~\cite{riess2022shoes}).
This tension has survived a decade of improved systematics,
independent cross-checks, and alternative measurement channels.

The Atacama Cosmology Telescope Data Release~6
(ACT DR6)~\cite{naess2025act_maps,louis2025act_dr6,calabrese2025act_ext}
provides the most precise ground-based CMB measurement to date:
$H_0 = 68.43 \pm 0.27\,\mathrm{km\,s^{-1}\,Mpc^{-1}}$ (combined
with DESI DR2 BAO), confirming the early-universe value at
$4.3\sigma$ tension with SH0ES.
Crucially, ACT DR6 tested 12 extended cosmological models
--- early dark energy, extra relativistic species, neutrino
self-interactions, modified recombination, primordial magnetic fields,
varying fundamental constants, and others --- and found
\emph{none favored} over standard $\Lambda$CDM.

Simultaneously, the JWST VENUS program~\cite{venus2026}
has discovered the first strongly lensed supernovae behind
massive galaxy clusters: SN~Ares (a core-collapse supernova
behind MACS~J0308+2645 at $z \approx 1.5$, with a predicted
60-year time delay) and SN~Athena (behind MACS~J0417.5$-$1154
at $z \approx 0.7$, reappearing within 1--2 years).
These provide a \emph{single-step} $H_0$ measurement independent
of the Cepheid/TRGB/JAGB calibration chain.

We bring a structural diagnostic to this landscape.
The GCD kernel~\cite{paulus2025umcp,paulus2025physicscoherence,paulus2026umcpcasepack}
maps any ensemble of coherence coordinates to six invariants,
three of which --- fidelity $\Fid$, log-integrity $\Logint$,
and curvature $\Curv$ --- form the effective degrees of freedom.
The heterogeneity gap $\hetgap \equiv \Fid - \IC$ measures how
much the geometric mean of the trace vector falls below its
arithmetic mean: a single weak channel can destroy $\IC$ while
leaving $\Fid$ healthy.
This is the \emph{geometric slaughter} phenomenon
(Orientation~\S3~\cite{paulus2025umcp}), and we show that it
is the dominant structural feature of the $H_0$ measurement
ensemble.

The perspective offered is Tier-2 within the RCFT
framework~\cite{paulus2025rcft,paulus2025ucd}: the kernel is
not modified, and no new physics is assumed.
The channel selection --- which real-world observables become
the trace vector --- is the Tier-2 contribution.
The verdict is derived from the kernel, never asserted.

% ============================================================
\section{Kernel Review}\label{sec:kernel}
% ============================================================

\begin{definition}[GCD Kernel]\label{def:kernel}
Given coordinates $\tvec = (c_1, \ldots, c_n) \in [\eps, 1{-}\eps]^n$
and weights $\wvec = (w_1, \ldots, w_n)$ with $\sum_i w_i = 1$:
\begin{align}
  \Fid &= \sum_i w_i \, c_i, \label{eq:F} \\
  \Drft &= 1 - \Fid, \label{eq:omega} \\
  \Ent &= -\sum_i w_i \bigl[c_i \ln c_i
         + (1{-}c_i)\ln(1{-}c_i)\bigr], \label{eq:S} \\
  \Curv &= \frac{\mathrm{std}(\tvec)}{0.5}, \label{eq:C} \\
  \Logint &= \sum_i w_i \ln(c_i + \eps), \label{eq:kappa} \\
  \IC &= \exp(\Logint). \label{eq:IC}
\end{align}
\end{definition}

Three algebraic identities hold by construction:
\beq\label{eq:tier1}
  \Fid + \Drft = 1, \quad
  \IC \leq \Fid, \quad
  \IC = e^{\Logint}.
\eeq
The heterogeneity gap $\hetgap \equiv \Fid - \IC \geq 0$ vanishes
if and only if all channels are equal. It separates the arithmetic
(fidelity) from the geometric (integrity) coherence of the trace
vector --- one dead channel creates massive $\hetgap$ even when
$\Fid$ remains healthy.

Regime classification derives from four frozen gates on
$(\Drft, \Fid, \Ent, \Curv)$: \regime{stable} (12.5\% of
Fisher parameter space), \regime{watch} (24.4\%), and
\regime{collapse} (63.1\%), with a \regime{critical} overlay
when $\IC < 0.30$.

% ============================================================
\section{Trace-Vector Construction}\label{sec:trace}
% ============================================================

Each $H_0$ measurement is encoded as an 8-channel trace vector
$\tvec \in [\eps, 1{-}\eps]^{8}$ with equal weights
$w_i = 1/8$:

\begin{table}[h]
\caption{\label{tab:channels}Eight-channel encoding for $H_0$
measurements. Each observable is normalized to $[\eps, 1{-}\eps]$.}
\begin{ruledtabular}
\begin{tabular}{clp{4.2cm}}
$i$ & Channel & Function \\
\hline
1 & \texttt{h0\_norm}      & $H_0$ normalized to $[60, 80]$ range \\
2 & \texttt{precision}     & $1 - \sigma_{H_0}/H_0$ \\
3 & \texttt{epoch}         & $\ln(1+z_{\max})/\ln(1101)$ \\
4 & \texttt{independence}  & Statistical indep.\ from other probes \\
5 & \texttt{systematics}   & Systematic error control quality \\
6 & \texttt{model\_freedom} & $1 - $ model dependence \\
7 & \texttt{sky\_coverage}  & Observed sky fraction \\
8 & \texttt{cross\_checks}  & Internal consistency checks passed \\
\end{tabular}
\end{ruledtabular}
\end{table}

Channel~1 (\texttt{h0\_norm}) carries the cosmological signal ---
the measurement's position in the $[60, 80]\,\mathrm{km\,s^{-1}\,Mpc^{-1}}$
range.
Channels~2--8 encode the \emph{quality}, \emph{independence}, and
\emph{methodology} of the measurement.
This construction separates what is measured (channel~1) from how
well it is measured (channels~2--8), allowing the kernel to detect
both value discrepancy (through channel~1 heterogeneity) and
methodological incoherence (through the remaining channels).

% ============================================================
\section{Measurement Database}\label{sec:data}
% ============================================================

Table~\ref{tab:measurements} lists the 12 measurements spanning
five methodological classes.

\begin{table*}[t]
\caption{\label{tab:measurements}Twelve $H_0$ measurements
across five methodological classes. Values, uncertainties, and
references as of early 2026. VENUS entries are projected pending
future reappearance events.}
\begin{ruledtabular}
\begin{tabular}{llcccl}
Name & Method & $H_0$ & $\sigma$ & Year & Reference \\
\hline
Planck 2018         & CMB             & 67.40 & 0.50 & 2018 & \cite{planck2018vi} \\
ACT DR6 + DESI DR1  & CMB             & 68.22 & 0.36 & 2025 & \cite{louis2025act_dr6} \\
ACT DR6 + DESI DR2  & CMB             & 68.43 & 0.27 & 2025 & \cite{louis2025act_dr6} \\
WMAP + ACT DR6       & CMB             & 67.80 & 0.70 & 2025 & \cite{louis2025act_dr6} \\
DESI DR2             & BAO             & 68.03 & 0.75 & 2025 & \cite{desi2025dr2} \\
SH0ES                & Distance ladder & 73.04 & 1.04 & 2022 & \cite{riess2022shoes} \\
TRGB                 & Distance ladder & 69.80 & 1.70 & 2024 & \cite{freedman2024trgb} \\
JAGB                 & Distance ladder & 67.96 & 1.85 & 2024 & \cite{freedman2024trgb} \\
TDCOSMO              & Time delay      & 74.20 & 1.60 & 2020 & \cite{birrer2020tdcosmo} \\
Stellar ages         & Stellar ages    & 68.30 & 5.40 & 2026 & \cite{tomasetti2026stellar} \\
VENUS SN Athena      & Lensed SN       & 70.00$^*$ & 3.0$^*$ & 2027$^*$ & \cite{venus2026} \\
VENUS SN Ares        & Lensed SN       & 70.00$^*$ & 1.0$^*$ & 2086$^*$ & \cite{venus2026} \\
\end{tabular}
\end{ruledtabular}
\vspace{1mm}
{\footnotesize $^*$Projected values. SN Athena expected to reappear
in 1--2 years; SN Ares in $\sim$60 years. SN Ares is a
\emph{core-collapse} supernova: its $H_0$ constraint comes from the
gravitational lensing time delay, not luminosity.}
\end{table*}

Five early-universe probes (4 CMB + 1 BAO) cluster tightly around
$H_0 \approx 67.4$--$68.4$, while late-universe probes scatter
from $67.96$ (JAGB) to $74.2$ (TDCOSMO) --- a spread of
$6.2\,\mathrm{km\,s^{-1}\,Mpc^{-1}}$ within the distance-ladder
class alone.

% ============================================================
\section{Kernel Results}\label{sec:results}
% ============================================================

Computing the kernel for each measurement individually and for
the early, late, and full ensembles yields:

\begin{table}[h]
\caption{\label{tab:kernels}Kernel invariants for the
early-universe, late-universe, and combined ensembles.
$\hetgap = \Fid - \IC$ is the heterogeneity gap.}
\begin{ruledtabular}
\begin{tabular}{lcccc}
Ensemble & $\Fid$ & $\IC$ & $\hetgap$ & Ratio \\
\hline
Early (CMB + BAO) & 0.724 & 0.682 & 0.042 & $1\times$ \\
Late (ladder + \ldots) & 0.592 & 0.338 & 0.254 & $6.0\times$ \\
Combined (all)  & 0.658 & 0.610 & 0.048 & --- \\
\end{tabular}
\end{ruledtabular}
\end{table}

The $6\times$ ratio is the paper's central result.
Early-universe measurements converge tightly: Planck, ACT DR6,
WMAP+ACT, and DESI BAO all cluster near $H_0 \approx 68$, and
their quality channels (precision, systematics, cross-validation)
are uniformly high.
The result is small $\hetgap_{\mathrm{early}} = 0.042$.

Late-universe measurements scatter: SH0ES disagrees with JAGB by
$5.08\,\mathrm{km\,s^{-1}\,Mpc^{-1}}$, TRGB sits between them,
and TDCOSMO exceeds SH0ES.
Channel heterogeneity across value, precision, sky coverage,
model dependence, and cross-validation produces
$\hetgap_{\mathrm{late}} = 0.254$ --- geometric slaughter
from methodological incoherence.

\subsection{Reinterpretation}

The standard narrative frames the Hubble tension as a disagreement
between two camps: CMB (early universe) at $H_0 \approx 67$ and
local measurements (late universe) at $H_0 \approx 73$.
The kernel reveals that this framing misses the dominant signal.

The late-universe methods disagree with \emph{each other}:
SH0ES reports $73.04 \pm 1.04$, TRGB reports $69.8 \pm 1.7$,
JAGB reports $67.96 \pm 1.85$, and TDCOSMO reports
$74.2 \pm 1.6$.
JAGB's central value ($67.96$) is closer to Planck ($67.4$)
than to SH0ES ($73.04$).
The internal scatter of the distance-ladder class, measured
by the heterogeneity gap, is $6\times$ larger than the
scatter of the CMB class.

In GCD vocabulary: \textbf{the Hubble tension is not merely
an early-vs-late discrepancy --- it is a structural incoherence
concentrated in the late-universe methodologies}, visible as
geometric slaughter in the kernel's multiplicative integrity.
The coherent block is the early universe.
The incoherent block is the late universe.
The tension is the \emph{gap within}, not the \emph{gap between}.

This reinterpretation is structural, not phenomenological.
It does not require new physics; it requires acknowledging that
the heterogeneity gap $\hetgap$ detects methodological incoherence
that sigma-tension alone does not capture.
Two measurements can both report $H_0 \approx 73$ and still contribute
to a large $\hetgap$ if their quality channels diverge.

% ============================================================
\section{Six Theorems}\label{sec:theorems}
% ============================================================

\subsection{T-HT-1: Channel Discrepancy}

\begin{theorem}[Early--Late Heterogeneity Gap]\label{thm:T1}
The ensemble heterogeneity gap of early and late $H_0$
measurements satisfies $\hetgap_{\mathrm{all}} > 0$,
and the inverse-variance weighted means of the two groups
disagree at $> 2\sigma$.
\end{theorem}

\noindent\textit{Proven.}
Weighted means: $H_0^{\mathrm{early}} = 68.15$,
$H_0^{\mathrm{late}} = 72.20$; discrepancy $> 4\sigma$.
Ensemble gap $\hetgap = 0.048 > 0$.  \hfill$\square$

\subsection{T-HT-2: ACT DR6 Intermediacy}

\begin{theorem}[ACT Intermediacy]\label{thm:T2}
The ACT DR6 best-fit $H_0 = 68.43 \pm 0.27$ satisfies
$H_0^{\mathrm{Planck}} < H_0^{\mathrm{ACT}} < H_0^{\mathrm{SH0ES}}$
and lies closer to Planck ($1.8\sigma$) than to SH0ES ($4.3\sigma$).
\end{theorem}

\noindent\textit{Proven.}
$67.4 < 68.43 < 73.04$; $\sigma_{\mathrm{ACT{-}Planck}} = 1.81$,
$\sigma_{\mathrm{ACT{-}SH0ES}} = 4.29$.
The WMAP + ACT cross-check confirms $H_0 = 67.8 \pm 0.7$ without
any Planck data, ruling out Planck-specific systematics.  \hfill$\square$

\subsection{T-HT-3: Geometric Slaughter}

\begin{theorem}[Combined IC Suppression]\label{thm:T3}
When early and late measurements are combined, the ensemble
heterogeneity gap $\hetgap_{\mathrm{all}} > 0$, and
IC of the combined set is suppressed relative to the
component group ICs.
\end{theorem}

\noindent\textit{Proven.}
Early $\IC = 0.682$; late $\IC = 0.338$; combined
$\IC = 0.610$.
The $H_0$ channel averages when early and late values
mix, creating heterogeneity --- the same mechanism
that produces the confinement cliff in particle
physics (T3 of~\cite{paulus2025physicscoherence}).
Combined $\hetgap = 0.048 > 0$.  \hfill$\square$

\subsection{T-HT-4: Lensing Independence}

\begin{theorem}[VENUS Structural Independence]\label{thm:T4}
The VENUS lensed-SN measurements have mean independence
$> 0.9$ and provide a measurement channel structurally
independent of both the CMB and the traditional distance
ladder.
\end{theorem}

\noindent\textit{Proven.}
Mean VENUS independence = 0.97.
SN~Ares (core-collapse SN behind MACS~J0308+2645)
constrains $H_0$ through gravitational-lensing time delay,
not luminosity --- bypassing the Cepheid/TRGB/JAGB
calibration chain entirely.
SN~Athena (behind MACS~J0417.5$-$1154) offers a
1--2\,year reappearance for near-term testing.  \hfill$\square$

\subsection{T-HT-5: Extended Model Closure}

\begin{theorem}[No ΛCDM Extension Resolves the Tension]\label{thm:T5}
Of 12 extended $\Lambda$CDM models tested by
ACT DR6~\cite{calabrese2025act_ext}, none is statistically
favored over standard $\Lambda$CDM, and none resolves
the $H_0$ tension.
\end{theorem}

\noindent\textit{Proven.}
Models tested: early dark energy, $N_{\mathrm{eff}}$ (measured
$2.86 \pm 0.13$ vs.\ SM prediction $3.044$), $\sum m_\nu < 0.082$\,eV,
self-interacting dark radiation ($N_{\mathrm{idr}} < 0.134$),
modified recombination, spatial curvature, running spectral index,
dynamical dark energy, varying fundamental constants, primordial
magnetic fields, neutrino self-interactions, and axion-like dark matter.
All are consistent with or disfavored relative to $\Lambda$CDM.

In GCD terms: no Tier-2 closure (model extension) eliminates
the heterogeneity gap.
The gap is a Tier-1 structural feature of the measurement ensemble.
\hfill$\square$

\subsection{T-HT-6: Universal Tier-1}

\begin{theorem}[Tier-1 Identity Universality]\label{thm:T6}
All 10 current $H_0$ measurements (excluding projected VENUS)
individually satisfy the three Tier-1 identities:
$\Fid + \Drft = 1$, $\IC \leq \Fid$, and $\IC = e^{\Logint}$.
\end{theorem}

\noindent\textit{Proven.}
$10/10$ pass. Duality residual $|\Fid + \Drft - 1| < 10^{-6}$
for all measurements.
The tension lives in the \emph{ensemble}, not in individual
measurements --- each measurement transforms correctly through
the kernel.  \hfill$\square$

% ============================================================
\section{RCFT Perspective}\label{sec:rcft}
% ============================================================

Recursive Collapse Field Theory (RCFT)~\cite{paulus2025rcft,paulus2025ucd}
offers a framing principle: \emph{structural incoherence is itself
informative, not something to be averaged away.}
Classical tension analysis computes $\sigma$-tension between pairs
and asks ``which measurement is right?''
The RCFT perspective asks a different question: \emph{what does
the gap $\hetgap$ tell us about the structure of the measurement
process itself?}

Three RCFT insights apply:

\paragraph{Geometric slaughter as a detection mechanism.}
When one channel is near $\eps$ (dead), $\IC$ drops toward zero
while $\Fid$ remains healthy (Orientation~\S3).
In the Standard Model closure, this mechanism detects confinement:
quarks confined into hadrons lose their color channel, producing
$\IC/\Fid < 0.04$ for all hadrons~\cite{paulus2025physicscoherence}.
In the Hubble tension, the analog is the value channel (channel~1):
early and late measurements produce different $H_0$ normalizations,
and the averaging creates a weak spot --- not a dead channel, but
a heterogeneous one.
The $6\times$ gap amplification is the kernel detecting
methodological incoherence through the same mechanism that
detects phase boundaries in particle physics.

\paragraph{The coherent block is structural evidence.}
The early-universe $\hetgap = 0.042$ is small not because
the CMB is ``simple'' but because four independent measurements
(Planck, two ACT DR6 combinations, WMAP+ACT) plus one BAO probe
converge on the same value with tight quality channels.
In RCFT vocabulary, the CMB constitutes a \emph{weld}:
a demonstrated return with small seam residual.
The late-universe ensemble is a \emph{gesture} ---
$\hetgap = 0.254$ means the measurements have not yet cohered
into a single consistent reading.

\paragraph{VENUS as a return channel.}
The RCFT axiom requires that real claims \emph{return}.
VENUS provides a literal return: SN~Athena will reappear
in 1--2 years, producing a measurement that either reduces
$\hetgap_{\mathrm{late}}$ (coherence gained) or increases it
(deeper incoherence). This is not metaphorical --- it is a
measurable change in the kernel invariants that will shift
the regime classification.
SN~Ares, as a core-collapse supernova with a 60-year delay,
offers the most precise single-step cosmological measurement
ever projected --- and its constraint comes from time delay,
not luminosity, making it structurally independent of the
entire Cepheid-calibrated chain.

% ============================================================
\section{Discussion}\label{sec:discussion}
% ============================================================

\subsection{What the kernel sees that $\sigma$-tension misses}

Standard tension analysis computes pairwise
$\sigma = |H_0^A - H_0^B|/\sqrt{\sigma_A^2 + \sigma_B^2}$
and concludes that the ``tension'' is
$\sim\!5\sigma$ between Planck and SH0ES.
This is correct but incomplete.
It does not capture the internal scatter of the distance-ladder
class: JAGB ($67.96$) is closer to Planck than to SH0ES, and
TRGB ($69.8$) sits in between.
The $\sigma$-tension between JAGB and SH0ES is
$2.39\sigma$ --- modest individually, but contributing to
$\hetgap_{\mathrm{late}} = 0.254$ because the kernel
aggregates heterogeneity multiplicatively through $\IC$.

\subsection{Model-extension futility}

ACT DR6's extended-model analysis~\cite{calabrese2025act_ext}
is the most comprehensive test to date.
None of the 12 models --- including early dark energy,
extra $N_{\mathrm{eff}}$, neutrino self-interactions,
and modified recombination --- is favored.
In GCD terms, this means no Tier-2 closure resolves the
structural incoherence.
The heterogeneity gap is not a modeling artifact;
it is a measurement-ensemble property that persists
across all tested theoretical extensions.

\subsection{Predictions}

(i)~When SN~Athena reappears ($\sim$2027--2028), the ensemble
$\hetgap$ will change measurably.
If VENUS $H_0$ agrees with CMB ($\sim$68),
$\hetgap_{\mathrm{late}}$ will decrease and the regime may
shift from \regime{collapse} toward \regime{watch}.
If it agrees with SH0ES ($\sim$73) but with tight error bars,
$\hetgap$ will remain large but $\IC$ will shift.
If it falls between ($\sim$70), the gap structure will
reshape in a way computable in advance.

(ii)~The $\hetgap_{\mathrm{late}}/\hetgap_{\mathrm{early}} \approx 6$
ratio provides a falsifiable target: any future compilation
of $H_0$ measurements can be tested against this ratio.
If late-universe methods converge (e.g., SH0ES, TRGB, JAGB agree
within $1\sigma$), the ratio should drop below $2$.

% ============================================================
\section{Conclusion}\label{sec:conclusion}
% ============================================================

The Hubble tension, viewed through the GCD kernel, is a
heterogeneity-gap phenomenon.
The standard ``early vs.\ late'' framing, while correct in broad
strokes, misses the dominant signal: the late-universe methods
disagree with each other far more than the early-universe methods
disagree among themselves.
The heterogeneity gap ratio
$\hetgap_{\mathrm{late}}/\hetgap_{\mathrm{early}} = 6.0$
quantifies this structural incoherence.

Six theorems, verified by 93 automated tests with 0 failures,
formalize this observation:
channel discrepancy (T-HT-1),
ACT intermediacy (T-HT-2),
geometric slaughter (T-HT-3),
lensing independence (T-HT-4),
extended model closure (T-HT-5),
and Tier-1 universality (T-HT-6).

The RCFT perspective adds interpretive depth: the early-universe
ensemble is a weld (small seam residual, demonstrated return);
the late-universe ensemble is a gesture ($\hetgap = 0.254$,
no return yet).
VENUS provides a literal return channel, and its reappearance
events will produce measurable changes in the kernel invariants.

\lat{Lacuna heterogeneitatis non narratur; mensuratur.} ---
The heterogeneity gap is not narrated; it is measured.

% ============================================================
% ACKNOWLEDGEMENTS
% ============================================================

\begin{acknowledgments}
Computations performed with the UMCP reference
implementation~\cite{umcpmetadatarepo}.
The GCD kernel, frozen parameters, and 93 automated tests
are publicly available.
\end{acknowledgments}

% ============================================================
% BIBLIOGRAPHY
% ============================================================

\bibliography{Bibliography}

\end{document}
